\section{Pairs and Lists}
\label{sec:lists}

\subsection{Pairs}

\code{cons} glues any two values into a \emph{pair}, and \code{car} and \code{cdr} take the pair apart:
\begin{lstlisting}[style=repl]
> (cons 3 4)
'(3 . 4)
> (car (cons 1 2))
1
> (cdr (cons 1 2))
2
\end{lstlisting}
The names \code{car} and \code{cdr} are historical, from the registers of the IBM~704 on which Lisp was first implemented. \code{first} and \code{rest} are the same operations on lists, with readable names (the names \cite{HtDP} uses). Both work; pick one style and stay with it.

\subsection{Lists Are Nested Pairs}

A list is not a separate kind of thing. It is a chain of pairs in which the right component is always another list, ending in \code{empty} (also written \code{'()}):

\begin{definition}[List]
A list is either \code{empty}, or \rkt|(cons x l)| where \code{x} is any value and \code{l} is a list.
\end{definition}

\begin{lstlisting}
(define l (cons 1 (cons 2 (cons 3 (cons 4 empty)))))
(define l-short (list 1 2 3 4))
\end{lstlisting}
\code{l} and \code{l-short} are the same value; \code{list} is shorthand for the chain of \code{cons}es (\Cref{fig:boxptr}). Both print as \code{'(1 2 3 4)}. Then \rkt|(first l)| is \code{1} and \rkt|(rest l)| is \code{'(2 3 4)}.

\begin{figure}[htbp]
\centering
\begin{tikzpicture}[x=1cm, y=1cm]
  \foreach \i/\v in {0/1,1/2,2/3,3/4} {
    \node[cell] (a\i) at (\i*2.1,0) {};
    \node[cell, right=0pt of a\i] (d\i) {};
    \node[cell, below=6mm of a\i, draw=none] (v\i) {\v};
    \draw[ptr] (a\i.center) -- (v\i.north);
  }
  \foreach \i/\j in {0/1,1/2,2/3} \draw[ptr] (d\i.center) -- (a\j.west);
  \draw (d3.south west) -- (d3.north east);
  \node[right=2mm of d3, lbl] {\code{empty}};
  \node[above=1mm of a0, lbl] {\code{car}};
  \node[above=1mm of d0, lbl] {\code{cdr}};

  \node[cell] (p) at (0,-2.4) {};
  \node[cell, right=0pt of p] (q) {};
  \node[below=6mm of p] (p3) {\code{3}};
  \node[below=6mm of q] (q4) {\code{4}};
  \draw[ptr] (p.center) -- (p3.north);
  \draw[ptr] (q.center) -- (q4.north);
  \node[right=4mm of q, lbl, align=left] {\rkt|(cons 3 4)| = \code{'(3 . 4)}:\\ the right side is not a list};
\end{tikzpicture}
\caption{Box-and-pointer diagrams. Top: \code{'(1 2 3 4)}, four pairs whose \code{cdr}s chain to \code{empty} (drawn as a slash). Bottom: a bare pair.}
\label{fig:boxptr}
\end{figure}

\begin{pitfall}[The dot]
The dot in \code{'(3 . 4)} means the right side of that pair is not a list, so the chain does not end in \code{empty}. Whenever a dot appears in output where you expected a list, you built a bare pair somewhere.
\end{pitfall}

The two-case definition is why list functions keep having two \code{cond} clauses: one for \code{empty}, one for a \code{cons}.

\subsection{Reversing a List}

\begin{keypoint}[The shape of the function follows the shape of the data]
A list is empty or a \code{cons}, so a function on lists has those two cases. In the \code{cons} case, the recursive call is on the \code{rest}, the part of the data that is itself a list.
\end{keypoint}

The empty list reversed is the empty list. Otherwise, reverse everything after the first element, then put the first element at the end:

\begin{lstlisting}
(define (my-reverse l)
  (cond
    [(empty? l) empty]
    [else (append (my-reverse (rest l)) (list (first l)))]))
\end{lstlisting}

\code{append} joins two lists, so both arguments must be lists. That is why the first element is wrapped as \rkt|(list (first l))|, the one-element list. Rewriting by hand:
\begin{lstlisting}[style=repl]
(my-reverse '(1 2 3))
---> (append (my-reverse '(2 3)) '(1))
---> (append (append (my-reverse '(3)) '(2)) '(1))
---> (append (append '(3) '(2)) '(1))
---> (append '(3 2) '(1))
---> '(3 2 1)
\end{lstlisting}

This works, but \code{append} itself walks its whole left argument, and it is called once per element, so \code{my-reverse} takes $O(n^2)$ time. Racket's built-in \code{reverse} uses an accumulator instead, the same trick as \code{fact-tail}; Recitation~1 writes that version (\Cref{sec:rec1}).
